\chapter{Discusión}
\label{cap:discusion}

El Capítulo~\ref{cap:introduccion} planteó dos preguntas, y los resultados las responden con distinto grado
de firmeza. La primera se refería a la calidad de las imágenes del algoritmo propuesto. Esas imágenes
fueron mejores que las del Top-Hat clásico en cuatro de las seis métricas. Frente a los cinco métodos
multiescala, la calidad del algoritmo fue comparable, aunque con un perfil propio: más actividad y menos
fidelidad. Sus imágenes tienen más información, contraste y nitidez, pero se parecen menos a las
fuentes. La propuesta encabezó la clasificación por rango promedio, pero por un margen estrecho y sin
que ese margen se sometiera a prueba. La segunda pregunta era si un detector entrenado con esas imágenes
encuentra más objetos vistos por una sola cámara que uno entrenado con las del Top-Hat clásico. El
detector encontró algo más, pero la ventaja es mínima. La tasa equilibrada, que promedia el acierto en
los objetos de solo la visible y de solo la infrarroja, fue 3.0 puntos mayor que la del Top-Hat
clásico, con un intervalo de 0.0 a 6.0. Según el par de semillas, esa diferencia cambia de signo. Así,
en calidad la hipótesis se sostiene en lo principal; en la detección, la ventaja queda como un hallazgo
por confirmar. El capítulo discute esos
resultados en cinco ejes: el perfil de calidad y su costo, el ajuste de los parámetros, la aplicación a
la detección, la relación entre calidad y detección y, al final, las limitaciones del estudio.

\section{La calidad de la fusión y su costo}
\label{sec:disc-calidad}

El primer puesto por rango promedio, con 3.642 frente a 3.783 de la DTCWT (Tabla~\ref{tab:ranking_seis}),
no autoriza a decir que la propuesta supere a ese método. Esa diferencia no tiene una prueba propia, y los
dos métodos llegan a la cabeza por caminos opuestos. La propuesta fue primera en actividad y sexta en
fidelidad; la DTCWT, sexta en actividad y segunda en fidelidad. Como cada grupo reúne tres métricas, el
rango global lo decide la suma de los dos rangos parciales: 7.28 para la propuesta y 7.57 para la DTCWT.
Un perfil marcado no basta, por sí solo, para llegar a la cabeza. La DWT (db4), todavía más
desequilibrada que la DTCWT entre los dos grupos, quedó quinta.

El perfil de la propuesta nace de la regla de fusión (Ecuación~\ref{eq:regla-tophat}), que toma su forma del realce de
contraste de la Ecuación~\eqref{eq:realce} \parencite{bai2012}. Sumar el detalle claro y restar el oscuro
agranda la diferencia de brillo entre cada detalle y su entorno, pero aleja la imagen de sus fuentes. El
barrido del peso lo confirma (Sección~\ref{sec:res-ajuste}). Con la propuesta y con el Top-Hat clásico,
subir $m$ elevó la desviación estándar y el gradiente medio y bajó las tres métricas de fidelidad. El
compromiso es, entonces, propio de la regla y no depende del banco de elementos. Algo parecido observaron
\textcite{bajacfigueredo2024} con otra variable. Al agrandar el elemento estructurante de 3 a 9,
subieron la desviación estándar, la entropía y la frecuencia espacial, y la información mutua bajó o se
mantuvo.

El costo se ve también en las imágenes. La resta del detalle oscuro es el único término de la regla que
lleva un píxel por debajo del promedio de las fuentes, y puede bajarlo hasta el negro, donde el valor se
recorta. De ahí el borde oscuro que rodea a la persona en la Figura~\ref{fig:montaje_kaptein}.

Frente al Top-Hat clásico, la propuesta ganó en entropía, desviación estándar, información mutua y
similitud estructural, perdió en gradiente medio y empató en PSNR (Tabla~\ref{tab:wilcoxon_resumen}). Ese
reparto no puede atribuirse al banco, porque entre los dos métodos cambian a la vez los elementos, el radio
y el peso. Según el barrido, el peso mueve por sí solo tres de esas métricas en la dirección
observada. Con los dos operadores, un peso menor reduce el gradiente medio y eleva la
información mutua y la similitud estructural. La propuesta usa un peso de 0.30, y el Top-Hat clásico,
de 1, así que parte de esas diferencias puede deberse solo a ese peso menor.

El primer puesto depende también del conjunto de métricas. Con las nueve métricas iniciales, cinco de
ellas de actividad, la propuesta quedó primera. Con un conjunto de diecisiete, en el que doce comparan
la fusión con sus fuentes, bajó al tercer puesto, detrás de la pirámide de Laplace y de la DTCWT. Por
otra parte, los cinco métodos multiescala se usaron con sus parámetros por defecto, mientras que el radio
de la propuesta se eligió sobre los mismos veinte pares de la evaluación (Sección~\ref{sec:ajuste}).
Nada asegura que el orden se mantenga si todos se ajustan con un mismo criterio.

El control con imágenes degradadas marca el límite de las métricas de actividad: una ventaja en ellas
no prueba, por sí sola, que el detalle agregado sea útil. Cuanto más ruido se agregó al promedio simple,
más puestos ganó la imagen resultante, hasta empatar en el octavo con el Top-Hat clásico
(Sección~\ref{sec:res-robustez}). Con la mayor intensidad fue, además, la mejor de las catorce entradas
en las tres métricas de actividad. En ese mismo control, la propuesta quedó tercera, detrás de la
pirámide de Laplace y del promedio simple sin degradar. Ese puesto lo producen las dos versiones
desenfocadas, que en fidelidad quedaron delante de la propuesta y, en las seis métricas, detrás del promedio
simple. Al ordenar solo los siete métodos y el promedio, la propuesta quedó primera, y el promedio, sexto. La cabeza de la clasificación resistió también el análisis por
escenas. Con un bloque por escena,
la propuesta siguió primera y la DTCWT segunda, y solo uno de los sesenta y seis contrastes cambió su
decisión.

\section{El ajuste de los parámetros}
\label{sec:disc-ajuste}

Las veinticinco corridas del enjambre terminaron con el peso en 0.30, el límite inferior del rango de
búsqueda. La grilla de comprobación indica que ese valor es un borde y no un máximo interior. En cada
radio, el mejor peso quedó muy por debajo de ese límite. A partir de 0.30, la aptitud, la función que
el enjambre busca maximizar, bajó con cada aumento del peso (Sección~\ref{sec:res-ajuste}). Una
explicación posible está en la forma de la aptitud (Ecuación~\ref{eq:aptitud}). Dos de sus tres términos, la similitud estructural y la PSNR, premian el
parecido con las fuentes, y cada aumento del peso aleja de ellas la fusión. La entropía, el tercer
término, no compensa esa pérdida.

El resultado difiere del de \textcite{ortegarodriguez2025}, que buscaron el peso en el mismo rango.
Hicieron 125 corridas, con veinticinco configuraciones en cada una de cinco escenas, y el peso terminó
en 0.30 solo en 17, todas de una misma escena. Aquel trabajo usó la aptitud original, que aquí se
adaptó, ajustó los parámetros escena por escena y empleó un disco único. Cualquiera de esas diferencias
podría explicar el contraste, que no se investigó. La última tiene una lectura directa. El detalle del
banco puede llegar al doble del de un disco del mismo radio (Ecuación~\ref{eq:banco}), así que un
mismo peso aleja más la fusión de sus fuentes.

En el radio, el enjambre no llegó a un valor único, porque con $m = 0.30$ la aptitud tiene dos máximos
locales: uno en $r = 1$ y otro, menor, en $r = 25$. Casi todas las corridas, incluidas las quinientas de
control, terminaron en uno de esos dos radios (Sección~\ref{sec:res-ajuste}), y la mayor aptitud
corresponde a $r = 1$. Aun así, se adoptó $r = 25$ porque prioriza el realce: con $r = 1$, el disco y
los segmentos miden apenas tres píxeles de ancho y solo extraen detalles muy finos. Frente a $r = 1$, el
radio adoptado fue mejor en todas las métricas de actividad y peor en todas las de fidelidad. Esa
elección explica en parte la ventaja de la propuesta sobre los métodos multiescala en entropía,
desviación estándar y gradiente medio.

\section{La aplicación a la detección}
\label{sec:disc-deteccion}

La ventaja de la propuesta frente al Top-Hat clásico, de 3.0 puntos de tasa equilibrada, admite dos
lecturas (Tabla~\ref{tab:deteccion_tophat}). La primera se apoya en el intervalo de confianza del
95~\%, de 0.0 a 6.0, que excluye el cero por muy poco: su límite inferior es de 0.02 puntos. El cero
queda fuera también con los otros dos umbrales, pero no con el \textit{bootstrap} salvaje por grupos,
más confiable cuando los grupos son pocos y de tamaño desigual. Con el umbral principal, ese cálculo
dio un intervalo de $-3.5$ a 5.5 y un valor de $p$ de 0.17. La segunda lectura surge de los pares de
semillas. Sus nueve diferencias van de $-3.1$ a 7.8 puntos, y las tres
negativas enfrentan a la propuesta con una misma semilla del Top-Hat clásico.

Las dos lecturas no se contradicen. El intervalo mide cómo varía la diferencia entre los grupos de
cuadros. Como la tasa de cada objeto llega ya promediada sobre las tres semillas, el intervalo no recoge
por completo la variación del entrenamiento. Los pares de semillas, en cambio, la dejan a la vista: basta que un
detector del Top-Hat clásico acierte más que los otros dos para invertir el signo en tres de los nueve
pares. Por eso la ventaja es mínima, y queda como un hallazgo por confirmar con más grupos, más semillas
y otros detectores (Capítulo~\ref{cap:recomendaciones}).

El hallazgo, además, no aísla el efecto del banco. Entre la propuesta y el Top-Hat clásico cambian a la vez
los elementos, el radio y el peso (Sección~\ref{sec:operador}), y la comparación no dice cuál de esos
cambios produce la diferencia. Tampoco muestra que el detalle agregado ayude a detectar. El promedio
simple es la regla de la Ecuación~\eqref{eq:regla-tophat} con peso cero, de modo que la comparación con
él mide el aporte del detalle. El Top-Hat clásico tuvo una tasa menor que la de ese promedio
(58.4~\% contra 60.2~\%, sin prueba entre ambos). Como referencia, la propuesta superó a ese promedio
por apenas 1.2 puntos. El intervalo de esa diferencia, calculado del mismo modo que el de la
Tabla~\ref{tab:deteccion_tophat}, va de $-6.6$ a 9.0 e incluye el cero.

Aun si se confirmara, la ventaja frente al Top-Hat clásico no pondría a la propuesta entre las mejores
entradas para detectar. Entre las ocho entradas fusionadas quedó cuarta, con 61.4~\%, y la pirámide de Laplace
tuvo la mayor tasa, 66.5~\% (Tabla~\ref{tab:v2_tasas}). Esa pirámide superó a la propuesta con todas las
combinaciones de semillas. En la mAP50, una medida secundaria calculada sobre todas las clases, cinco de
las ocho entradas fusionadas superaron a la infrarroja sola. La propuesta no estuvo entre ellas, aunque
su diferencia con la infrarroja fue menor que la variación entre sus propias semillas
(Sección~\ref{sec:res-secundarias}). En esa medida y en los cuadros con un objeto de cada cámara,
además, el Top-Hat clásico quedó apenas por encima de la propuesta, sin prueba entre ambos.

Estas comparaciones valen para un entrenamiento de cuarenta épocas. En el seguimiento hecho con la
primera semilla de cada entrada, nueve de los diez detectores todavía mejoraban al final del
entrenamiento su mAP50 sobre los cuadros de seguimiento, que no son los de la prueba. Entre las épocas 30 y 40, además, la propuesta pasó de estar por debajo a estar por encima
de la pirámide de Laplace y de la DWT (db4) en esa medida (Sección~\ref{sec:res-secundarias}). Con un
entrenamiento más largo, el orden entre entradas de tasa parecida podría cambiar.

\section{Relación entre calidad y detección}
\label{sec:disc-relacion}

Los dos planos de la evaluación no ordenaron a los métodos de la misma manera. La correlación de Spearman
entre los dos órdenes fue de $-0.32$ y no resultó significativa (Sección~\ref{sec:res-secundarias}).
Con siete métodos la prueba es débil, y solo una relación muy fuerte habría alcanzado la significación.
Los dos órdenes provienen, además, de conjuntos distintos: la calidad se midió en el TNO y la detección
en M3FD. Los dos métodos morfológicos ilustran la distancia entre la actividad y la detección. Fueron
los mejores en actividad, pero en detección, entre las ocho entradas fusionadas, la propuesta quedó
cuarta y el Top-Hat clásico, séptimo. Esto no prueba que la calidad sea irrelevante para detectar, pero
el orden por métricas de imagen no anticipó el orden en la tarea. Ese desacuerdo respalda evaluar la
fusión también por su efecto en una tarea, como proponen \textcite{liu2022} y \textcite{tang2022}.

\section{Limitaciones del estudio}
\label{sec:disc-limitaciones}

Los resultados permiten precisar las limitaciones anunciadas en la Sección~\ref{sec:limitaciones}. En
calidad, el primer puesto no equivale a superioridad. Descansa en diferencias de rango promedio que no se
sometieron a prueba y cambia con el conjunto de métricas. Ese puesto se ganó en actividad a costa de la
fidelidad. Además, se obtuvo en condiciones que pueden favorecer a la propuesta: su radio se eligió
sobre los mismos veinte pares de la evaluación, y los métodos con los que se la comparó no se ajustaron.
Tres de esos veinte pares sirvieron también para calcular la aptitud del enjambre. Todo lo
dicho vale para una sola configuración, $r = 25$ y $m = 0.30$, y el barrido muestra que el perfil se
desplaza con el peso. El radio se mide en píxeles, sin adaptarse al tamaño de la imagen. En el TNO, el disco de 51 píxeles de ancho
cubre entre el 8.9~\% y el 20.6~\% del lado menor de la imagen. Un mismo radio abarca, así, detalles
de distinto tamaño relativo.

En la detección, los objetos de una sola cámara quedaron definidos por los detectores de la visible y de
la infrarroja, sin revisión humana. Los objetos de ese tipo se reparten en solo diez grupos de cuadros, y
con tan pocos grupos la comparación tiene poca potencia, es decir, poca capacidad de revelar una diferencia
real. Cuando los grupos son pocos y de tamaño muy desigual, además, el error estándar robusto por grupos tiende a
quedarse corto, y el intervalo puede resultar más estrecho de lo debido. Sin el mayor de esos grupos,
la diferencia con el Top-Hat clásico baja de 3.0 a 1.5 puntos. Las luces, por su parte, están
concentradas: casi dos de cada tres provienen de un solo grupo (Sección~\ref{sec:caracterizacion}). Se usó un solo tipo de detector, de tamaño pequeño, con tres
semillas por entrada y cuarenta épocas de entrenamiento.

El estudio no incluyó métodos de fusión basados en redes neuronales ni los morfológicos multiescala de
\textcite{bai2013} y \textcite{roman2020}. El foco estuvo en un operador de una sola escala, y la
comparación con esas familias queda para trabajos futuros (Capítulo~\ref{cap:recomendaciones}).
